% SHARD-ID: CA-67-REFINEMENT-REQUIREMENTS-OBSTRUCTIONS
% SHARD-TITLE: Refinement-Map Requirements and Obstructions
% SHARD-SUMMARY: Collects the sourced OAR conditions a refinement family must satisfy, and the soft relaxations available when exact intertwining fails.
% SHARD-SUMMARY: Extracts the mechanism of the Jones and Kliesch-Koenig no-go results as a checklist of failure modes to avoid.
% SHARD-SUMMARY: Locates precisely where fixed anyon number is baked into the braiding-RG refinement, motivating the variable-N family of CA-68.
% SHARD-KEYWORDS: refinement maps, scaling maps, OAR, soft inductive limit, Jones no-go, Kliesch-Koenig, braiding RG, Koo-Saleur

\section{Refinement-Map Requirements and Obstructions}
\label{sec:refinement-requirements-obstructions}

\subsection{Status, Sources, and Dependencies}

\textbf{Status: source-local synthesis plus derived checklist.}  Every
requirement and failure mode below is anchored in a registered local source;
the checklist assembly and the dense-chain diagnosis are local readings.

Dependencies: CA-14, CA-15, CA-62, CA-65, CA-66; CONVENTIONS.md (r).

% Source: literature/md/2010.11121/2010.11121.md:101--111 -- scaling maps are
%   unital, completely positive, *-morphisms, with the semigroup property (2.6).
% Source: literature/md/2010.11121/2010.11121.md:170 -- directed systems of
%   C*-algebras use unital injective *-morphisms; injective implies isometric.
% Source: literature/md/2010.11121/2010.11121.md:206--210 -- bounded
%   N-independent support growth; the scaling map identifies a lattice field
%   with a smeared continuum field, not a point-like one.
% Source: literature/md/2010.11121/2010.11121.md:156--164 -- dynamics
%   convergence: refined evolutions must be Cauchy in the scaling-limit GNS
%   representation; sensitivity to the choice of scaling maps.
% Source: literature/md/2010.11121/2010.11121.md:1372--1374 -- block-spin fails
%   the algebra-level limit, yet the vacuum is recovered via Wightman two-point
%   functions (Theorem 5.1).
% Source: literature/md/2010.11121/2010.11121.md:1444 -- point-localization:
%   continuum local operators only via Wightman-type reconstruction.
% Source: literature/md/2306.16063/2306.16063.md:185 -- soft inductive systems:
%   linear contractions with asymptotic transitivity.
% Source: literature/md/2306.16063/2306.16063.md:355--358 -- operations
%   converge iff they asymptotically intertwine the connecting maps (eq. 21).
% Source: literature/md/2306.16063/2306.16063.md:571--575 -- soft C*-systems:
%   asymptotic multiplicativity condition (eq. 29).
% Source: literature/md/2306.16063/2306.16063.md:1025--1057 -- Thompson-group
%   actions exist iff connecting maps come from a single local refinement;
%   implementers fail for nonzero central charge; Koo-Saleur succeeds.
% Source: references/lattice-symmetry/Jones2017NoGo/source/nogo.tex:61--69 --
%   no-go abstract: block-spin/TL semicontinuous space carries no chiral CFT
%   with rotations as limits of lattice rotations.
% Source: references/lattice-symmetry/Jones2017NoGo/source/nogo.tex:592--594 --
%   root cause: the dangling binary tree does not feel the circle topology.
% Source: references/lattice-symmetry/Jones2017NoGo/source/nogo.tex:694--704 --
%   the renormalisation (squaring) map is norm-contracting; iterates vanish.
% Source: references/lattice-symmetry/Jones2017NoGo/source/nogo.tex:951--955 --
%   main theorem: triadic-rotation-rigid vectors vanish.
% Source: references/lattice-symmetry/KlieschKoenig2020/source/ms.tex:329 --
%   discontinuous representations have full measure among isometries; an easily
%   verified necessary condition for a continuous limit.
% Source: references/lattice-symmetry/KlieschKoenig2020/source/ms.tex:420--432 --
%   the necessary-condition lemma implying the main theorem.
% Source: literature/md/2201.11562/2201.11562.md:48--50 -- anyonic-chain basis:
%   fusion-tree labels with a sigma-anyon at every site.
% Source: literature/md/2201.11562/2201.11562.md:70--74 -- the refining step
%   attaches through-strings and braids their ends into position, increasing
%   the local density of sigma-anyons; Reidemeister moves give morphism property.
% Source: references/text/OARWavelets.txt:393--398 -- continuous translations
%   recovered from dyadic ones via the wavelet scaling map.
% Source: references/text/OARWavelets.txt:405--406 -- Lorentz/conformal
%   convergence requires further work.

\subsection{Strict OAR Requirements}

The sourced strict architecture asks of a refinement family
\(\alpha_{N'}^{N}:\Alg_N\to\Alg_{N'}\):
\[
\begin{array}{ll}
1.&\text{unital and completely positive (states pull back to states);}\\
2.&\text{a *-morphism, injective, hence isometric;}\\
3.&\text{the semigroup property }
   \alpha_{N''}^{N'}\circ\alpha_{N'}^{N}=\alpha_{N''}^{N};\\
4.&\text{bounded, scale-independent support growth (locality);}\\
5.&\text{smearing: a lattice field maps to a smeared continuum field,}\\
  &\text{not a point-like one.}
\end{array}
\]
On the state side, a scaling limit is a projectively consistent family
\(\omega^{(N)}=\omega^{(N')}\circ\alpha_{N'}^{N}\); on the dynamics side, the
refined evolutions must form Cauchy sequences in the scaling-limit GNS
representation, and the source stresses that success depends sensitively on
the choice of scaling maps.

\subsection{Fallback and Soft Relaxations}

Two sourced fallbacks weaken the strict demands without abandoning the
continuum target.  First, when the *-morphism property or regularity fails
(block-spin, point-localization), the algebra-level limit is lost but the
vacuum can still be recovered through Wightman two-point functions; continuum
local operators are then only indirectly accessible by reconstruction.

Second, the soft-inductive-system formalism replaces exact structure by
asymptotic structure: connecting maps need only be contractions with
asymptotic transitivity; operations converge iff they \emph{asymptotically
intertwine} the connecting maps,
\[
  \lim_{n\gg m}\bigl\|(\tilde\jmath_{nm}T_m-T_n j_{nm})x_m\bigr\|=0,
\]
and C*-structure survives under an asymptotic multiplicativity condition.
For symmetry recovery the same source is explicit: Thompson-group implementers
fail for nonzero central charge, and the successful strategy is Koo--Saleur
approximation of conformal symmetries.  \textbf{Consequence [derived]:} for a
CFT target with \(c\neq0\), ``compatible with the Virasoro action'' can only
mean \emph{asymptotic} intertwining of normal-ordered Koo--Saleur generators
on a suitable core --- there is no exact-implementer formulation to aim for.

\subsection{The No-Go Mechanism as a Checklist}

The Jones no-go concerns the semicontinuous limit built from a single
scale-invariant tree-local refinement element.  The registered source
identifies the root cause: the underlying dangling binary tree does not feel
the circle topology, and the pulled-back smallest rotation is a long product
of local elements controlled by a norm-contracting renormalisation (squaring)
map, whose iterates vanish --- forcing rotation matrix elements weakly to
zero.  Kliesch--Koenig extend this: discontinuity has \emph{full measure}
among tree isometries, and they isolate an easily verified necessary
condition for a continuous limit.

\textbf{Checklist [derived reading].}  A refinement family risks the no-go
mechanism when:
\[
\begin{array}{ll}
\mathrm{NG1}&\text{it is generated by a single scale-invariant tree-local}\\
            &\text{element, decoupled from the circle;}\\
\mathrm{NG2}&\text{its induced squaring/renormalisation map is norm-}\\
            &\text{contracting on the relevant subspace;}\\
\mathrm{NG3}&\text{it embeds point-like, without low-pass smearing;}\\
\mathrm{NG4}&\text{it operates in the large-loop-parameter regime where}\\
            &\text{the contraction estimate bites;}\\
\mathrm{NG5}&\text{it demands exact unitary implementers of the discrete}\\
            &\text{rotation/conformal action on a fixed limit space.}
\end{array}
\]
Any candidate family must either avoid each mode or verify the
Kliesch--Koenig necessary condition explicitly.  The wavelet precedent shows
what success looks like: translation-covariant smearing recovers continuous
translations, while Lorentz/conformal covariance is separate further work.

\subsection{Where Fixed Anyon Number Enters the Braiding RG}

The braiding-RG refinement for anyonic chains is sourced as: attach new
through-strings to a string-box diagram and braid their ends into position,
\emph{increasing the local density of \(\sigma\)-anyons}; the basis carries a
\(\sigma\)-anyon at every site.  \textbf{Diagnosis [derived].}  Three
fixed-\(N\) commitments follow: (i) the braid needs a strand at every site it
acts on, so the map is undefined on empty sites; (ii) the identification of
local operators with Temperley--Lieb generators presupposes exactly one anyon
per site; (iii) the refinement doubles the strand count deterministically.
For the CA-65/CA-66 site object \(O=\one\oplus X\), occupation is a quantum
variable, so a refinement must be defined on the whole word space including
empty sites --- the braiding-RG step does not port.  Conversely, dense chains
have no refinable vacuum: every refined site must carry an anyon, so any
fine-graining changes the physical content.  The variable-\(N\) tower is
exactly what removes this obstruction, at the price of the new proof
obligations recorded in CA-68.

\subsection{Boundary}

This shard proves no convergence theorem and constructs no refinement; it
fixes the requirement set and the failure modes.  The Kliesch--Koenig
necessary condition is registered but its verification for any variable-\(N\)
family is future work.  The Koo--Saleur compatibility target is stated
precisely in CA-68; nothing here claims it holds.
